\scanpage{80}
% Archive pencil foliation: 79.
valente à $\handwritten{\underline}{G}(\handwritten{\ell})$ pour tout corps alg. clos $\handwritten{\ell}$ de ca\rlap{\raisebox{-.65em}{r}}r. $p\handwritten{>0}$ (de façon précise, si $\handwritten{\overline}{\handwritten{\underline}{F}}_p$ est\deleted{\unsure{l|t}} la clôture algébrique de $\handwritten{\underline}{F}_p$ dans $\handwritten{\ell}$, l'inclusion de celle-ci dans $\handwritten{\ell}$ ind\unsure{\deleted{i}u|\deleted{u}i}it une équivalence de \rlap{\raisebox{.4em}{g}}gerbes
\[
\handwritten{\underline}{G}(\handwritten{\ell})\handwritten{\underline{\approx}}\handwritten{\underline}{G}(\handwritten{\overline}{\handwritten{\underline}{F}}_p)\deleted{\handwritten{\unsure{l|)|,|;}}}\handwritten{\underline{\approx}}\handwritten{\underline}{H}\deleted{\handwritten{\unsure{)|l|.}}}\qquad ).
\]
8.6. Etude de $\operatorname{Rep}(\handwritten{\underline}{H})$. Nous avons dit que les classes d'objets simples sont en correspondance biunivoque avec les élé\unsure{me|\deleted{m}e|\deleted{e}m} nt s de $\handwritten{\underline}{Q}$ ; l'élément de $\handwritten{\underline}{Q}$ correspondant à un objet isotypique est \unsure{\deleted{e}d|\deleted{d}e|d}it la \handwritten{\uline}{pente} de cet élément. Nous avons vu que l'objet de $\operatorname{Rep}(\handwritten{\underline}{H})$ défini \margin{\handwritten{Notation $E_{s/r}$ pour $E_{r,s}$ indiquée, \unsure{donc|dans}
\[
E_\lambda\otimes E_{\lambda'}\simeq E_{\lambda,\lambda'}^{\unsure{n|m}(\lambda,\lambda')}
\]
pour $\lambda,\lambda'\in\mathbb{Q}$}}
par l'objet $\handwritten{\underline}{Q}_p[T]/(T^r-p^s)$ de $\operatorname{Fcriso}(\handwritten{\underline}{F}_p)$, muni de la multiplication par $T$, est simple de pente $s/r$ ($s,r\handwritten{\in}\handwritten{\underline}{Z}$, $r\geq1$, $(s,r)=1$ ), de sorte que nous obtenons ainsi un catalogue complet des objets simples $E_{r,s}$. Bien entendu, $E_{r,s}$ est de rang $r$, d'autre part son corps des endomorphismes est \deleted{\unsure{lta|lla|ila}} un corps gauche \inserted{$K_{r,s}$}\handwritten{\begin{tikzpicture}[baseline=0cm,x=1cm,y=1cm]\draw[thin] (0,0) .. controls (.35,.16) and (.7,.1) .. (1.05,.08) .. controls (1.45,.05) and (1.6,.22) .. (1.64,.48);\end{tikzpicture}} dont la classe dans $\operatorname{Br}(\handwritten{\underline}{Q}_p)=\handwritten{\underline}{Q}/\handwritten{\underline}{Z}$ est $s/r\bmod1$, d'après ce qu'on a vu dans 8.3.

$\operatorname{Rep}(\handwritten{\underline}{H})$ est limite inductive des sous-c\unsure{\deleted{q}a|\deleted{a}q|a}tégories pleines $\operatorname{Rep}(\handwritten{\underline}{H},r)$, $r\geq1$, où $(\handwritten{\underline}{H},r)$ est le quotient de la gerbe $\handwritten{\underline}{H}$ correspond au quotient $(H,r)\handwritten{\simeq}$\handwritten{\begin{tikzpicture}[baseline=-.1cm,x=1cm,y=1cm]\draw[gray,thin] (.02,.15)..controls (.04,-.22) and (.23,-.23)..(.28,-.08);\end{tikzpicture}}$\handwritten{\underline}{G}_m$ de son lien, tel que l'homomorphisme de Tate \deleted{soit un} induise un homomorphisme \begin{tikzcd}[cramped](H,r)\arrow[r,"\handwritten{\unsure{q|a|\alpha}}"]&\handwritten{\underline}{G}_m\end{tikzcd} s'identifiant à \begin{tikzcd}[cramped]x\arrow[r,mapsto]&x^r\end{tikzcd} : F-iso\inserted{cristaux (sur $\handwritten{\overline}{k}$) de pente multiple de $1/r$.}\handwritten{\begin{tikzpicture}[baseline=0cm,x=1cm,y=1cm]\draw[thin] (0,-.17)..controls (.16,.05) and (.35,.02)..(.62,.02)..controls (.85,.03) and (.95,.08)..(1.05,.12);\end{tikzpicture}} Donc $(\handwritten{\underline}{H},r)$ \deleted{s'identifie} \unsure{\deleted{a}è|\deleted{è}a|è}s\unsure{\deleted{l}t|\deleted{t}l|t}\deleted{\unsure{a|à}} \inserted{une}\handwritten{\begin{tikzpicture}[baseline=0cm,x=1cm,y=1cm]\draw[thin] (0,-.13)--(.65,.13);\end{tikzpicture}} gerbe de lien $\handwritten{\underline}{G}_m$, ayant comme classe dans $\operatorname{Br}(\handwritten{\underline}{Q}_p)$ l\unsure{\deleted{s}e|\deleted{e}s|\deleted{i}e|\deleted{e}i|e} générateur canonique $1/r\bmod\handwritten{\underline}{Z}$ de ${}_r\operatorname{Br}(\handwritten{\underline}{Q}_p)\handwritten{\simeq}\handwritten{\underline}{Z}/r\handwritten{\underline}{Z}$. C\rlap{,}\unsure{'|I}est aussi, plus canoniquement, ``la'' gerbe définie par le corps gauche $K_{r,\handwritten{1}\deleted{\unsure{s|a|8}}}$ plus haut, \deleted{qui}

\handwritten{[}Exercice amusant, que je n'ai pas fait: la structure additive de $\operatorname{Rep}(\handwritten{\underline}{H})$ étant entièrement décrite \unsure{pa|\deleted{p}a|\deleted{a}p} r la collection de\unsure{s|\deleted{z}s|\deleted{s}z} $K_{r,s}$, expliciter sa $\handwritten{\otimes}$-structure ... i.e. écrire les \handwritten{``}relations canoniques\handwritten{''} en termes de prod\unsure{\deleted{n}u|\deleted{u}n|\deleted{l}u|\deleted{u}l|u|n}i\unsure{\deleted{s}t|\deleted{t}s|\deleted{e}t|\deleted{t}e|t|s|e}s tensoriels entre les $E_{r,s}$. Autre exerciee suggéré: trouver ``canoniquement'' une description des modules de Dieudonné des réduits des groupes de Lubin-Tate en termes des $E_{r,s}$ \deleted{\handwritten{\unsure{;|,|)|.}}} \handwritten{\unsure{(|[}}en particulier, les Lubin-Tate associés aux extensions non ramifiées

